% Appendix A
\section{Программын кодын хэсгээс}
\label{app:code}

\subsection{Захиалга хүлээн авах функц (Atomic Transaction)}
Энэхүү функц нь хүргэгч захиалгыг хүлээн авах үед өгөгдлийн бүрэн бүтэн байдлыг хангах зорилгоор Firestore Transaction ашигладаг.

\begin{lstlisting}[language=JavaScript, caption=Accept Order Transaction]
export const acceptOrder = async (orderId: string, courierId: string) => {
  try {
    await runTransaction(db, async (transaction) => {
      // 1. Get order ref
      const orderRef = doc(db, "orders", orderId);
      const orderDoc = await transaction.get(orderRef);

      if (!orderDoc.exists()) {
        throw "Order does not exist!";
      }

      // 2. Get courier info for denormalization
      const courierRef = doc(db, "users", courierId);
      const courierDoc = await transaction.get(courierRef);
      
      if (!courierDoc.exists()) {
        throw "Courier profile not found!";
      }

      const courierData = courierDoc.data();
      
      // 3. Prepare denormalized data
      const courierInfo: CourierInfo = {
        name: courierData.name || courierData.displayName || "Courier",
        phone: courierData.phone || courierData.phoneNumber || "",
      };
      
      if (courierData.photoURL) {
        courierInfo.photoURL = courierData.photoURL;
      }

      // 4. Update order atomically
      transaction.update(orderRef, {
        status: "accepted",
        courierId: courierId,
        acceptedAt: serverTimestamp(),
        updatedAt: serverTimestamp(),
        courierInfo: courierInfo 
      });
    });
    
    console.log("Order accepted successfully via transaction");
  } catch (error) {
    console.error("Transaction failed: ", error);
    throw error;
  }
};
\end{lstlisting}

\subsection{Firestore Security Rules}
Хэрэглэгчийн өгөгдлийг хамгаалах дүрмийн жишээ.

\begin{lstlisting}[caption=Firestore Security Rules]
rules_version = '2';
service cloud.firestore {
  match /databases/{database}/documents {
    
    // Helper functions
    function isAuthenticated() {
      return request.auth != null;
    }
    
    function isOwner(userId) {
      return request.auth.uid == userId;
    }

    // Users collection
    match /users/{userId} {
      allow read: if isAuthenticated();
      allow create: if isAuthenticated() && isOwner(userId);
      allow update: if isAuthenticated() && isOwner(userId);
    }
    
    // Orders collection
    match /orders/{orderId} {
      allow read: if isAuthenticated();
      allow create: if isAuthenticated();
      allow update: if isAuthenticated(); 
    }
  }
}
\end{lstlisting}

\subsection{Firebase тохиргоо (firebase.ts)}
\begin{lstlisting}[language=JavaScript, caption=Firebase Config]
import { initializeApp } from "firebase/app";
import { getAuth } from "firebase/auth";
import { getFirestore } from "firebase/firestore";
import { getStorage } from "firebase/storage";

const firebaseConfig = {
  apiKey: process.env.EXPO_PUBLIC_FIREBASE_API_KEY,
  authDomain: process.env.EXPO_PUBLIC_FIREBASE_AUTH_DOMAIN,
  projectId: process.env.EXPO_PUBLIC_FIREBASE_PROJECT_ID,
  storageBucket: process.env.EXPO_PUBLIC_FIREBASE_STORAGE_BUCKET,
  messagingSenderId: process.env.EXPO_PUBLIC_FIREBASE_MESSAGING_SENDER_ID,
  appId: process.env.EXPO_PUBLIC_FIREBASE_APP_ID,
};

const app = initializeApp(firebaseConfig);

export const auth = getAuth(app);
export const db = getFirestore(app);
export const storage = getStorage(app);

export default app;
\end{lstlisting}
